\section*{Exercises for Chapter~\ref{ch:prelim}}
\addcontentsline{toc}{section}{Exercises for Chapter 2}

\begin{enumerate}
  \item Prove that if $a\mid b$ and $b\mid c$, then $a\mid c$.
  \item Prove that if $a\mid b$ and $a\mid c$, then $a\mid (bx+cy)$ for all integers $x,y$.
  \item Use the Euclidean algorithm to compute $\gcd(1071,462)$. Write each step as $r_{k+1}=r_{k-1}-q_k r_k$.
  \item Find integers $x,y$ with $1071x+462y=\gcd(1071,462)$, by running the algorithm backwards.
  \item Show that $\gcd(a,b)=\gcd(a,b-ka)$ for every integer $k$.
  \item Prove that $\gcd(n,n+1)=1$ for every integer $n$.
  \item Prove that if $\gcd(a,b)=1$ and $\gcd(a,c)=1$, then $\gcd(a,bc)=1$.
  \item Let $p$ be prime. Prove that $p\mid a^2$ implies $p\mid a$.
  \item Find all integers $n$ with $1\le n\le 50$ for which $n$ and $n+2$ are both prime. These are the twin primes; state whether this list is expected to be infinite, and why this cannot be decided by the methods of this chapter.
  \item Prove by induction that $1+2+\dots+n=\frac{n(n+1)}{2}$ for every $n\ge 1$.
  \item Prove that $n^3-n$ is divisible by $6$ for every integer $n\ge 0$.
  \item Show that the division algorithm (Theorem~\ref{thm:division}) fails if $b=0$. Explain which step of the proof uses $b>0$.
\end{enumerate}
